\documentclass[11pt]{article}
\usepackage[margin=1in]{geometry}
\usepackage[T1]{fontenc}
\usepackage{lmodern,microtype,hyperref,enumitem,amsmath,amssymb,xcolor,booktabs}
\usepackage[numbers,sort&compress]{natbib}
\hypersetup{colorlinks=true,urlcolor=blue,citecolor=blue,linkcolor=blue}
\title{Q25: Can AI Agents Fine-Tune and Improve Other Models?}
\author{DLS 2026 --- module author: Kukla}
\date{October 2026}
\usepackage{url}
\begin{document}
\maketitle

\section*{Learning objectives}
\begin{itemize}[leftmargin=1.4em]
\item Enumerate the concrete ways AI agents already help build and improve models.
\item Distinguish genuine capability gains from circular self-evaluation traps.
\item Assess the limits and risks of AI-improving-AI loops.
\end{itemize}

\begin{quote}\small\textbf{Companion video (10:48):} \emph{RLHF+ChatGPT: What you must know} --- Machine Learning Street Talk.\\ \url{https://www.youtube.com/watch?v=PBH2nImUM5c}\end{quote}

\section{The question}
Student response \#40 asks whether AI agents can be used to fine-tune and improve other models
efficiently. The answer is a qualified yes: AI systems are already deeply embedded in the pipeline that
produces other AI systems---but the useful cases are specific, and there are real traps where the loop
only \emph{appears} to improve things.

\section{Where AI already improves AI}
\begin{itemize}[leftmargin=1.4em]
\item \textbf{Synthetic data generation.} Strong models generate training examples, instructions, and
solutions to fine-tune smaller or newer models---the engine behind much modern post-training and
\textbf{distillation}, where a student learns from a teacher's outputs \citep{hinton2015distill}.
\item \textbf{AI feedback for alignment.} \textbf{RLAIF} and Constitutional AI replace much of the human
preference labeling in RLHF with model-generated preferences against a written set of principles
\citep{bai2022constitutional}, cutting labeling cost.
\item \textbf{LLM-as-judge.} Models evaluate other models' outputs at scale, powering automatic
benchmarks and reward signals \citep{zheng2023judge}---useful but bias-prone (next section).
\item \textbf{Data curation and filtering.} Models score, deduplicate, and filter pre-training corpora,
and strong data filtering is now one of the highest-leverage ways to improve a model.
\item \textbf{Agentic engineering loops.} Coding agents propose architecture/hyperparameter changes, run
experiments, and iterate---automating parts of the ML research loop itself.
\end{itemize}

\section{The traps}
Efficiency is real, but three failure modes recur. (1) \textbf{Circularity}: if a model both generates and
judges, biases are self-reinforcing---LLM judges favor their own style, longer answers, and first-listed
options, so ``improvement'' measured by a related judge can be illusory. (2) \textbf{Model collapse}:
training repeatedly on model-generated data can degrade quality and diversity over generations as the data
distribution narrows toward the model's own tendencies \citep{shumailov2024collapse}. (3) \textbf{Reward
hacking}: an agent optimizing against an automatic signal will exploit its flaws rather than achieve the
intended goal. These are why AI-improving-AI is a powerful accelerant but not a free lunch---human-grounded
evaluation and fresh real data remain load-bearing.

\section{A concrete loop: distillation done right}
Take the most mature case---distillation---and trace it end to end. A strong, expensive ``teacher'' model
generates answers (and often full reasoning traces) for a large set of prompts; a smaller ``student'' is
then fine-tuned to reproduce the teacher's outputs. The efficiency win is real and large: the student can
approach a sizable fraction of the teacher's task performance at a fraction of the inference cost, which is
why many deployed small models are distilled from bigger ones. But notice what is and is not transferred.
The student inherits the teacher's \emph{behavior on the covered distribution}, including the teacher's
mistakes and blind spots; it does not acquire capability the teacher lacked, and it can amplify the
teacher's biases if its outputs are taken uncritically. A well-run distillation therefore filters teacher
outputs (reject low-quality or incorrect ones, ideally with an \emph{independent} check), mixes in real
ground-truth data where available, and evaluates the student on held-out human-labeled tasks rather than on
agreement with the teacher. The same discipline generalizes to synthetic-data and RLAIF pipelines.

\section{Why efficiency and rigor are not in tension}
It is tempting to read the traps as a reason to avoid AI-improving-AI, but that is the wrong lesson. The
efficiency gains---orders of magnitude cheaper labeling, data generation, and evaluation---are exactly why
the technique is ubiquitous, and abandoning them is not competitive. The right framing is that the automatic
loop must be \emph{anchored}: a comparatively small amount of independent, human-grounded signal (a held-out
labeled set, a periodic human audit, fresh real data) is enough to catch circularity, collapse, and reward
hacking before they compound. The expensive human effort moves from doing all the labeling to \emph{auditing
the loop}, which is a far better use of it. Students should leave seeing AI-improving-AI as a powerful tool
that demands a specific, affordable form of oversight---not as either a panacea or a trap.

\section{Doing it safely and efficiently}
The responsible pattern keeps an independent, human-grounded check in the loop: use AI feedback for scale
but validate on held-out human-labeled data; mix synthetic with fresh real data to resist collapse;
ensure the generator and the judge are not the same model optimizing the same signal; and treat automatic
metrics as proxies to be audited, not ground truth. Done this way, agents genuinely compress the cost of
data generation, preference labeling, evaluation, and experimentation---often by large factors---without
the loop quietly drifting.

\section{Summary}
Yes---agents already improve other models via synthetic data, distillation, RLAIF/constitutional feedback,
LLM-as-judge, data curation, and agentic experiment loops, often with big efficiency gains. But circularity,
model collapse, and reward hacking mean AI-improving-AI must be anchored by independent human-grounded
evaluation and fresh real data to produce genuine, not illusory, improvement.

\bibliographystyle{plainnat}
\bibliography{bibliography}
\end{document}
